\chapter{Introduction}

\section{What is Hayashi}

\hayashi{} is an interpreted programming language designed for applied
econometric analysis. The name honours Fumio Hayashi, whose book
\textit{Econometrics} (2000) is a canonical reference in the field.

The language has three central goals:

\begin{enumerate}
  \item \textbf{Familiar syntax} for Stata users — commands such as
        \lstinline{generate}, \lstinline{summarize}, and \lstinline{regress}
        work as expected.
  \item \textbf{General-purpose programming} — functions, closures, control
        structures, modules. It is not merely a collection of statistical
        commands.
  \item \textbf{Performance without configuration} — the interpreter is
        written in Rust; no JVM, no GC, no runtime dependencies.
\end{enumerate}

\section{Design philosophy}

\subsection*{Clear semantics before concise syntax}

When an operation mutates an object, it should look visibly imperative.
When it returns a new value, it should look functional. \hayashi{}
distinguishes these two forms consistently:

\begin{lstlisting}[caption={Imperative form vs.\ functional form}]
// imperative: modifies df in place (Stata-style)
generate df z = x^2

// functional: returns a new DataFrame, df unchanged
let df2 = generate(df, z = x^2)
\end{lstlisting}

\subsection*{Copy-on-Write for DataFrames}

DataFrames use copy-on-write (COW) semantics: assigning a DataFrame to
another variable creates a cheap alias. The actual copy only happens at
the moment of mutation. This allows efficient sharing without surprises:

\begin{lstlisting}[caption={COW in action}]
load "data.csv" as df
let baseline = df      // alias, no copy

let transformed = df |> generate(z = x^2) |> filter(z > 10)
// baseline still has no column z
\end{lstlisting}

\subsection*{Error messages as documentation}

Errors in \hayashi{} include the code snippet that caused them, the line
number, and, when possible, a suggested fix (\textit{did you mean?}).
A well-formatted error message is more useful than a manual.

\section{Hayashi and Stata}

\hayashi{} is not a Stata clone. The data command syntax is inspired by
Stata because that syntax is good — it is close to the language economists
use naturally. But \hayashi{} is a full language: it has first-class
functions, closures, recursion, and structured error handling.

The table below summarises the key differences:

\begin{center}
\begin{tabular}{lll}
\toprule
\textbf{Aspect} & \textbf{Stata} & \textbf{Hayashi} \\
\midrule
License         & Proprietary      & GPL-3.0 \\
Functions       & Limited          & First-class \\
Closures        & No               & Yes \\
Types           & Implicit         & Explicit in functions \\
DataFrames      & One at a time    & Multiple, COW \\
Modularisation  & \texttt{do}      & \lstinline{source} \\
Performance     & C                & Rust \\
\bottomrule
\end{tabular}
\end{center}

\section{Hayashi and R / Python}

R and Python have vast ecosystems. \hayashi{} does not compete with them
on the volume of available packages — it competes on the clarity of syntax
for the specific task of applied econometrics. A panel regression script in
\hayashi{} has less visual noise than its equivalent in any other language.

\section{Empirical validation programme}

\hayashi{} includes a reproducible validation programme in the
\texttt{validation/} directory. Each case runs a \hayashi{} script and
compares it, within declared tolerances, against reference implementations in
R and Python (statsmodels). The cases use both public real datasets and
simulated DGPs taken from the chapters of this book.

To run the full validation suite:

\begin{lstlisting}[caption={Run the validation programme}]
python -m venv validation/.venv
validation/.venv/bin/pip install -r validation/requirements.txt
hay validate
\end{lstlisting}

The current status of every case is recorded in \texttt{validation/MATRIX.md}.

\section{Structure of this book}

\begin{description}
  \item[Chapters 2--3] Installation and use of the REPL.
  \item[Chapters 4--12] Type system, variables, operators, control flow,
        functions, errors, strings, lists, and dictionaries.
  \item[Chapters 13--15] DataFrames, data manipulation, and the pipe
        operator.
  \item[Chapter 16] Modularisation with \lstinline{source}.
\end{description}

Parts II through VII cover descriptive statistics, OLS, IV, panel,
time series, limited dependent variable models, causal inference, and
advanced methods.
